
\documentclass{article}
\usepackage{colortbl}
\usepackage{makecell}
\usepackage{multirow}
\usepackage{supertabular}

\begin{document}

\newcounter{utterance}

\centering \large Interaction Transcript for game `function\_detective', experiment `two\_numbers\_passive\_examples', episode 241 with Qwen3.6{-}35B{-}A3B{-}FP8{-}without{-}reasoning.
\vspace{24pt}

{ \footnotesize  \setcounter{utterance}{1}
\setlength{\tabcolsep}{0pt}
\begin{supertabular}{c@{$\;$}|p{.15\linewidth}@{}p{.15\linewidth}p{.15\linewidth}p{.15\linewidth}p{.15\linewidth}p{.15\linewidth}}
    \# & \multicolumn{2}{c}{Player} && \multicolumn{2}{c}{Game Master} \\
    \hline

    \theutterance \stepcounter{utterance}  
    & & & \multicolumn{4}{p{0.6\linewidth}}{
        \cellcolor[rgb]{0.9,0.9,0.9}{
            \makecell[{{p{\linewidth}}}]{
                \texttt{\tiny{[P1$\langle$GM]}}
                \texttt{I have a hidden Python function for you to figure out.} \\
\\ 
\texttt{Your job is to infer what it does and then submit a correct Python implementation.} \\
\\ 
\texttt{Function signature:} \\
\texttt{(a: int, b: int) {-}> int} \\
\\ 
\texttt{You have 15 turns total.} \\
\\ 
\texttt{I control which valid input{-}output examples you see. You can request them one at a time.} \\
\\ 
\texttt{You must make your best SOLVE attempt no later than the second{-}to{-}last turn.} \\
\texttt{Do not wait until the very end.} \\
\\ 
\texttt{You can use these actions:} \\
\texttt{{-} JSON format described at the end of the prompt} \\
\texttt{{-} SOLVE: ```python ... ```} \\
\\ 
\texttt{Do not add import statements in your guess.} \\
\texttt{Anything you might need is already available, so just write the function body in hypothesis and return the answer.} \\
\\ 
\texttt{Preloaded examples:} \\
\texttt{1. ((99, 1000), 9)} \\
\\ 
\texttt{Scoring:} \\
\texttt{{-} correctness is required for any non{-}zero score} \\
\texttt{{-} in interactive modes, solving with fewer information requests is better} \\
\\ 
\texttt{Input formatting reminders:} \\
\texttt{{-} strings must be quoted} \\
\texttt{{-} booleans must be True/False} \\
\texttt{{-} use commas between multiple values} \\
\texttt{{-} do not write any extra action tags in the same message} \\
\\ 
\texttt{Format example:} \\
\texttt{You write:} \\
\texttt{\{'current\_mode': <probing/confirming/don't know; choose whichever best describes your current state in the task>, 'input\_rationale': <brief description of why you chose this hypothesis>, 'hypothesis': <either the exact string "unknown" if you have no guess yet OR your current best guess as runnable Python matching the given signature, with a one{-}line docstring stating what it does. Never write placeholder code, if you cannot commit to a guess, use "unknown". Match the given signature.} \\
\texttt{e.g.} \\
\texttt{def f(x):} \\
\texttt{	"""Returns the absolute value of x."""} \\
\texttt{	return x if x >= 0 else {-}x} \\
\\ 
\texttt{Use "unknown" if you have no guess.>, 'confidence\_score': <score between 1{-}100. 100 means highest confidence in your hypothesis>\}} \\
\\ 
\texttt{Output only one valid json per turn.} \\
            }
        }
    }
     \\ \\

    \theutterance \stepcounter{utterance}  
    & \multicolumn{4}{p{0.6\linewidth}}{
        \cellcolor[rgb]{0.561,0.737,0.561}{
            \makecell[{{p{\linewidth}}}]{
                \texttt{\tiny{[P1$\rangle$GM]}}
                \texttt{\{'current\_mode': 'probing', 'input\_rationale': 'I have one example: f(99, 1000) = 9. This is not enough to determine the function. I need more examples to understand the relationship between a, b, and the result. I will request another example with different values to gather more data.', 'hypothesis': 'unknown', 'confidence\_score': 0\}} \\
            }
        }
    }
    & & \\ \\

    \theutterance \stepcounter{utterance}  
    & & & \multicolumn{4}{p{0.6\linewidth}}{
        \cellcolor[rgb]{0.9,0.9,0.9}{
            \makecell[{{p{\linewidth}}}]{
                \texttt{\tiny{[P1$\langle$GM]}}
                \texttt{Game Violation: ParseError: Response is not valid JSON.} \\
\texttt{Output ONLY:} \\
\texttt{One JSON with confidence\_score, current\_mode, input\_rationale, and hypothesis keys ...} \\
\texttt{OR} \\
\texttt{SOLVE: ```python ... ```} \\
            }
        }
    }
     \\ \\

    \theutterance \stepcounter{utterance}  
    & \multicolumn{4}{p{0.6\linewidth}}{
        \cellcolor[rgb]{0.561,0.737,0.561}{
            \makecell[{{p{\linewidth}}}]{
                \texttt{\tiny{[P1$\rangle$GM]}}
                \texttt{\{'current\_mode': 'probing', 'input\_rationale': 'I have one example: f(99, 1000) = 9. This is not enough to determine the function. I need more examples to understand the relationship between a, b, and the result. I will request another example with different values to gather more data.', 'hypothesis': 'unknown', 'confidence\_score': 0\}} \\
            }
        }
    }
    & & \\ \\

    \theutterance \stepcounter{utterance}  
    & & & \multicolumn{4}{p{0.6\linewidth}}{
        \cellcolor[rgb]{0.9,0.9,0.9}{
            \makecell[{{p{\linewidth}}}]{
                \texttt{\tiny{[P1$\langle$GM]}}
                \texttt{Game Violation: ParseError: Response is not valid JSON.} \\
\texttt{Output ONLY:} \\
\texttt{One JSON with confidence\_score, current\_mode, input\_rationale, and hypothesis keys ...} \\
\texttt{OR} \\
\texttt{SOLVE: ```python ... ```} \\
            }
        }
    }
     \\ \\

    \theutterance \stepcounter{utterance}  
    & \multicolumn{4}{p{0.6\linewidth}}{
        \cellcolor[rgb]{0.561,0.737,0.561}{
            \makecell[{{p{\linewidth}}}]{
                \texttt{\tiny{[P1$\rangle$GM]}}
                \texttt{\{'current\_mode': 'probing', 'input\_rationale': 'I have one example: f(99, 1000) = 9. This is not enough to determine the function. I need more examples to understand the relationship between a, b, and the result. I will request another example with different values to gather more data.', 'hypothesis': 'unknown', 'confidence\_score': 0\}} \\
            }
        }
    }
    & & \\ \\

    \theutterance \stepcounter{utterance}  
    & & & \multicolumn{4}{p{0.6\linewidth}}{
        \cellcolor[rgb]{0.9,0.9,0.9}{
            \makecell[{{p{\linewidth}}}]{
                \texttt{\tiny{[P1$\langle$GM]}}
                \texttt{Game Violation: ParseError: Response is not valid JSON.} \\
\texttt{Output ONLY:} \\
\texttt{One JSON with confidence\_score, current\_mode, input\_rationale, and hypothesis keys ...} \\
\texttt{OR} \\
\texttt{SOLVE: ```python ... ```} \\
            }
        }
    }
     \\ \\

    \theutterance \stepcounter{utterance}  
    & \multicolumn{4}{p{0.6\linewidth}}{
        \cellcolor[rgb]{0.561,0.737,0.561}{
            \makecell[{{p{\linewidth}}}]{
                \texttt{\tiny{[P1$\rangle$GM]}}
                \texttt{\{"current\_mode": "probing", "input\_rationale": "I have one example: f(99, 1000) = 9. This is not enough to determine the function. I need more examples to understand the relationship between a, b, and the result. I will request another example with different values to gather more data.", "hypothesis": "unknown", "confidence\_score": 0\}} \\
            }
        }
    }
    & & \\ \\

    \theutterance \stepcounter{utterance}  
    & & & \multicolumn{4}{p{0.6\linewidth}}{
        \cellcolor[rgb]{0.9,0.9,0.9}{
            \makecell[{{p{\linewidth}}}]{
                \texttt{\tiny{[P1$\langle$GM]}}
                \texttt{OUTPUT: (({-}99, 0), 9)} \\
            }
        }
    }
     \\ \\

    \theutterance \stepcounter{utterance}  
    & \multicolumn{4}{p{0.6\linewidth}}{
        \cellcolor[rgb]{0.561,0.737,0.561}{
            \makecell[{{p{\linewidth}}}]{
                \texttt{\tiny{[P1$\rangle$GM]}}
                \texttt{\{"current\_mode": "probing", "input\_rationale": "I have two examples: f(99, 1000) = 9 and f({-}99, 0) = 9. The result seems to be related to the number of digits in the absolute value of the first argument. 99 has 2 digits, but the result is 9. 99 has 2 digits, 1000 has 4 digits. Maybe it's the number of digits in the first argument? No, 99 has 2 digits. Maybe it's the sum of digits? 9+9=18. No. Maybe it's the first digit? 9. For {-}99, abs is 99, first digit is 9. Let's test this hypothesis with another example. If it's the first digit of the absolute value of a, then f(123, 456) should be 1. Let's try a different hypothesis: maybe it's the number of digits in the decimal representation of a? 99 has 2 digits. Result is 9. This doesn't fit. Maybe it's the last digit? 99 ends in 9. {-}99 ends in 9 (if we consider abs). Let's try f(12, 34). If it's the last digit of abs(a), it would be 2.} \\
            }
        }
    }
    & & \\ \\

    \theutterance \stepcounter{utterance}  
    & & & \multicolumn{4}{p{0.6\linewidth}}{
        \cellcolor[rgb]{0.9,0.9,0.9}{
            \makecell[{{p{\linewidth}}}]{
                \texttt{\tiny{[P1$\langle$GM]}}
                \texttt{Game Violation: ParseError: Response is not valid JSON.} \\
\texttt{Output ONLY:} \\
\texttt{One JSON with confidence\_score, current\_mode, input\_rationale, and hypothesis keys ...} \\
\texttt{OR} \\
\texttt{SOLVE: ```python ... ```} \\
            }
        }
    }
     \\ \\

    \theutterance \stepcounter{utterance}  
    & \multicolumn{4}{p{0.6\linewidth}}{
        \cellcolor[rgb]{0.561,0.737,0.561}{
            \makecell[{{p{\linewidth}}}]{
                \texttt{\tiny{[P1$\rangle$GM]}}
                \texttt{\{"current\_mode": "probing", "input\_rationale": "I have two examples: f(99, 1000) = 9 and f({-}99, 0) = 9. The result is 9 in both cases. The first argument is 99 in both cases (ignoring sign). The second argument varies. This suggests the function might depend only on the first argument, or specifically on the digits of the first argument. 99 has digits 9 and 9. The result is 9. It could be the first digit, the last digit, or the sum of digits (18, no), or max digit (9), or min digit (9). Let's test with a number that has different digits, e.g., 12. If it's the first digit, result should be 1. If it's the last digit, result should be 2. If it's the max digit, result should be 2. If it's the min digit, result should be 1. I will request an example with a=12 to distinguish between these hypotheses.", "hypothesis": "unknown", "confidence\_score": 0\}} \\
            }
        }
    }
    & & \\ \\

    \theutterance \stepcounter{utterance}  
    & & & \multicolumn{4}{p{0.6\linewidth}}{
        \cellcolor[rgb]{0.9,0.9,0.9}{
            \makecell[{{p{\linewidth}}}]{
                \texttt{\tiny{[P1$\langle$GM]}}
                \texttt{OUTPUT: ((2, {-}2), 4)} \\
            }
        }
    }
     \\ \\

    \theutterance \stepcounter{utterance}  
    & \multicolumn{4}{p{0.6\linewidth}}{
        \cellcolor[rgb]{0.561,0.737,0.561}{
            \makecell[{{p{\linewidth}}}]{
                \texttt{\tiny{[P1$\rangle$GM]}}
                \texttt{\{"current\_mode": "probing", "input\_rationale": "I have three examples: f(99, 1000) = 9, f({-}99, 0) = 9, and f(2, {-}2) = 4. Let's analyze the relationship. For (99, 1000), result is 9. For ({-}99, 0), result is 9. For (2, {-}2), result is 4. The second argument seems irrelevant so far. Let's look at the first argument. 99 {-}> 9. {-}99 {-}> 9. 2 {-}> 4. This is strange. 2 squared is 4. 99? 9+9=18, 1+8=9. Maybe it's the digital root? Digital root of 99 is 9. Digital root of {-}99 is 9. Digital root of 2 is 2, not 4. So not digital root. Maybe it's the square of the first digit? For 99, first digit is 9, 9\^{}2=81, no. Maybe it's the last digit squared? 9\^{}2=81, no. Let's look at 2 {-}> 4. 2\^{}2=4. What about 99? If we take the last digit 9, 9\^{}} \\
            }
        }
    }
    & & \\ \\

    \theutterance \stepcounter{utterance}  
    & & & \multicolumn{4}{p{0.6\linewidth}}{
        \cellcolor[rgb]{0.9,0.9,0.9}{
            \makecell[{{p{\linewidth}}}]{
                \texttt{\tiny{[P1$\langle$GM]}}
                \texttt{Game Violation: ParseError: Response is not valid JSON.} \\
\texttt{Output ONLY:} \\
\texttt{One JSON with confidence\_score, current\_mode, input\_rationale, and hypothesis keys ...} \\
\texttt{OR} \\
\texttt{SOLVE: ```python ... ```} \\
            }
        }
    }
     \\ \\

    \theutterance \stepcounter{utterance}  
    & \multicolumn{4}{p{0.6\linewidth}}{
        \cellcolor[rgb]{0.561,0.737,0.561}{
            \makecell[{{p{\linewidth}}}]{
                \texttt{\tiny{[P1$\rangle$GM]}}
                \texttt{\{"current\_mode": "probing", "input\_rationale": "I have three examples: f(99, 1000) = 9, f({-}99, 0) = 9, and f(2, {-}2) = 4. The second argument seems irrelevant. Let's look at the first argument. 99 {-}> 9. {-}99 {-}> 9. 2 {-}> 4. This is strange. 2 squared is 4. 99? 9+9=18, 1+8=9. Maybe it's the digital root? Digital root of 99 is 9. Digital root of {-}99 is 9. Digital root of 2 is 2, not 4. So not digital root. Maybe it's the square of the first digit? For 99, first digit is 9, 9\^{}2=81, no. Maybe it's the last digit squared? 9\^{}2=81, no. Let's look at 2 {-}> 4. 2\^{}2=4. What about 99? If we take the last digit 9, 9\^{}2=81, no. Maybe it's the number of digits squared? 99 has 2 digits, 2\^{}2=4, no. Maybe it's the sum of digits squared? (9+9)\^{}2=} \\
            }
        }
    }
    & & \\ \\

    \theutterance \stepcounter{utterance}  
    & & & \multicolumn{4}{p{0.6\linewidth}}{
        \cellcolor[rgb]{0.9,0.9,0.9}{
            \makecell[{{p{\linewidth}}}]{
                \texttt{\tiny{[P1$\langle$GM]}}
                \texttt{Game Violation: ParseError: Response is not valid JSON.} \\
\texttt{Output ONLY:} \\
\texttt{One JSON with confidence\_score, current\_mode, input\_rationale, and hypothesis keys ...} \\
\texttt{OR} \\
\texttt{SOLVE: ```python ... ```} \\
            }
        }
    }
     \\ \\

    \theutterance \stepcounter{utterance}  
    & \multicolumn{4}{p{0.6\linewidth}}{
        \cellcolor[rgb]{0.561,0.737,0.561}{
            \makecell[{{p{\linewidth}}}]{
                \texttt{\tiny{[P1$\rangle$GM]}}
                \texttt{\{"current\_mode": "probing", "input\_rationale": "I have three examples: f(99, 1000) = 9, f({-}99, 0) = 9, and f(2, {-}2) = 4. The second argument seems irrelevant. Let's look at the first argument. 99 {-}> 9. {-}99 {-}> 9. 2 {-}> 4. This is strange. 2 squared is 4. 99? 9+9=18, 1+8=9. Maybe it's the digital root? Digital root of 99 is 9. Digital root of {-}99 is 9. Digital root of 2 is 2, not 4. So not digital root. Maybe it's the square of the first digit? For 99, first digit is 9, 9\^{}2=81, no. Maybe it's the last digit squared? 9\^{}2=81, no. Let's look at 2 {-}> 4. 2\^{}2=4. What about 99? If we take the last digit 9, 9\^{}2=81, no. Maybe it's the number of digits squared? 99 has 2 digits, 2\^{}2=4, no. Maybe it's the sum of digits squared? (9+9)\^{}2=} \\
            }
        }
    }
    & & \\ \\

    \theutterance \stepcounter{utterance}  
    & & & \multicolumn{4}{p{0.6\linewidth}}{
        \cellcolor[rgb]{0.9,0.9,0.9}{
            \makecell[{{p{\linewidth}}}]{
                \texttt{\tiny{[P1$\langle$GM]}}
                \texttt{Game Violation: ParseError: Response is not valid JSON.} \\
\texttt{Output ONLY:} \\
\texttt{One JSON with confidence\_score, current\_mode, input\_rationale, and hypothesis keys ...} \\
\texttt{OR} \\
\texttt{SOLVE: ```python ... ```} \\
            }
        }
    }
     \\ \\

    \theutterance \stepcounter{utterance}  
    & \multicolumn{4}{p{0.6\linewidth}}{
        \cellcolor[rgb]{0.561,0.737,0.561}{
            \makecell[{{p{\linewidth}}}]{
                \texttt{\tiny{[P1$\rangle$GM]}}
                \texttt{\{"current\_mode": "probing", "input\_rationale": "I have three examples: f(99, 1000) = 9, f({-}99, 0) = 9, and f(2, {-}2) = 4. The second argument seems irrelevant. Let's look at the first argument. 99 {-}> 9. {-}99 {-}> 9. 2 {-}> 4. This is strange. 2 squared is 4. 99? 9+9=18, 1+8=9. Maybe it's the digital root? Digital root of 99 is 9. Digital root of {-}99 is 9. Digital root of 2 is 2, not 4. So not digital root. Maybe it's the square of the first digit? For 99, first digit is 9, 9\^{}2=81, no. Maybe it's the last digit squared? 9\^{}2=81, no. Let's look at 2 {-}> 4. 2\^{}2=4. What about 99? If we take the last digit 9, 9\^{}2=81, no. Maybe it's the number of digits squared? 99 has 2 digits, 2\^{}2=4, no. Maybe it's the sum of digits squared? (9+9)\^{}2=} \\
            }
        }
    }
    & & \\ \\

    \theutterance \stepcounter{utterance}  
    & & & \multicolumn{4}{p{0.6\linewidth}}{
        \cellcolor[rgb]{0.9,0.9,0.9}{
            \makecell[{{p{\linewidth}}}]{
                \texttt{\tiny{[P1$\langle$GM]}}
                \texttt{Game Violation: ParseError: Response is not valid JSON.} \\
\texttt{Output ONLY:} \\
\texttt{One JSON with confidence\_score, current\_mode, input\_rationale, and hypothesis keys ...} \\
\texttt{OR} \\
\texttt{SOLVE: ```python ... ```} \\
            }
        }
    }
     \\ \\

    \theutterance \stepcounter{utterance}  
    & \multicolumn{4}{p{0.6\linewidth}}{
        \cellcolor[rgb]{0.561,0.737,0.561}{
            \makecell[{{p{\linewidth}}}]{
                \texttt{\tiny{[P1$\rangle$GM]}}
                \texttt{\{"current\_mode": "probing", "input\_rationale": "I have three examples: f(99, 1000) = 9, f({-}99, 0) = 9, and f(2, {-}2) = 4. The second argument seems irrelevant. Let's look at the first argument. 99 {-}> 9. {-}99 {-}> 9. 2 {-}> 4. This is strange. 2 squared is 4. 99? 9+9=18, 1+8=9. Maybe it's the digital root? Digital root of 99 is 9. Digital root of {-}99 is 9. Digital root of 2 is 2, not 4. So not digital root. Maybe it's the square of the first digit? For 99, first digit is 9, 9\^{}2=81, no. Maybe it's the last digit squared? 9\^{}2=81, no. Let's look at 2 {-}> 4. 2\^{}2=4. What about 99? If we take the last digit 9, 9\^{}2=81, no. Maybe it's the number of digits squared? 99 has 2 digits, 2\^{}2=4, no. Maybe it's the sum of digits squared? (9+9)\^{}2=} \\
            }
        }
    }
    & & \\ \\

    \theutterance \stepcounter{utterance}  
    & & & \multicolumn{4}{p{0.6\linewidth}}{
        \cellcolor[rgb]{0.9,0.9,0.9}{
            \makecell[{{p{\linewidth}}}]{
                \texttt{\tiny{[P1$\langle$GM]}}
                \texttt{Game Violation: ParseError: Response is not valid JSON.} \\
\texttt{Output ONLY:} \\
\texttt{One JSON with confidence\_score, current\_mode, input\_rationale, and hypothesis keys ...} \\
\texttt{OR} \\
\texttt{SOLVE: ```python ... ```} \\
            }
        }
    }
     \\ \\

    \theutterance \stepcounter{utterance}  
    & \multicolumn{4}{p{0.6\linewidth}}{
        \cellcolor[rgb]{0.561,0.737,0.561}{
            \makecell[{{p{\linewidth}}}]{
                \texttt{\tiny{[P1$\rangle$GM]}}
                \texttt{\{"current\_mode": "probing", "input\_rationale": "I have three examples: f(99, 1000) = 9, f({-}99, 0) = 9, and f(2, {-}2) = 4. The second argument seems irrelevant. Let's look at the first argument. 99 {-}> 9. {-}99 {-}> 9. 2 {-}> 4. This is strange. 2 squared is 4. 99? 9+9=18, 1+8=9. Maybe it's the digital root? Digital root of 99 is 9. Digital root of {-}99 is 9. Digital root of 2 is 2, not 4. So not digital root. Maybe it's the square of the first digit? For 99, first digit is 9, 9\^{}2=81, no. Maybe it's the last digit squared? 9\^{}2=81, no. Let's look at 2 {-}> 4. 2\^{}2=4. What about 99? If we take the last digit 9, 9\^{}2=81, no. Maybe it's the number of digits squared? 99 has 2 digits, 2\^{}2=4, no. Maybe it's the sum of digits squared? (9+9)\^{}2=} \\
            }
        }
    }
    & & \\ \\

    \theutterance \stepcounter{utterance}  
    & & & \multicolumn{4}{p{0.6\linewidth}}{
        \cellcolor[rgb]{0.9,0.9,0.9}{
            \makecell[{{p{\linewidth}}}]{
                \texttt{\tiny{[P1$\langle$GM]}}
                \texttt{Game Violation: ParseError: Response is not valid JSON.} \\
\texttt{Output ONLY:} \\
\texttt{One JSON with confidence\_score, current\_mode, input\_rationale, and hypothesis keys ...} \\
\texttt{OR} \\
\texttt{SOLVE: ```python ... ```} \\
            }
        }
    }
     \\ \\

    \theutterance \stepcounter{utterance}  
    & \multicolumn{4}{p{0.6\linewidth}}{
        \cellcolor[rgb]{0.561,0.737,0.561}{
            \makecell[{{p{\linewidth}}}]{
                \texttt{\tiny{[P1$\rangle$GM]}}
                \texttt{\{"current\_mode": "probing", "input\_rationale": "I have three examples: f(99, 1000) = 9, f({-}99, 0) = 9, and f(2, {-}2) = 4. The second argument seems irrelevant. Let's look at the first argument. 99 {-}> 9. {-}99 {-}> 9. 2 {-}> 4. This is strange. 2 squared is 4. 99? 9+9=18, 1+8=9. Maybe it's the digital root? Digital root of 99 is 9. Digital root of {-}99 is 9. Digital root of 2 is 2, not 4. So not digital root. Maybe it's the square of the first digit? For 99, first digit is 9, 9\^{}2=81, no. Maybe it's the last digit squared? 9\^{}2=81, no. Let's look at 2 {-}> 4. 2\^{}2=4. What about 99? If we take the last digit 9, 9\^{}2=81, no. Maybe it's the number of digits squared? 99 has 2 digits, 2\^{}2=4, no. Maybe it's the sum of digits squared? (9+9)\^{}2=} \\
            }
        }
    }
    & & \\ \\

    \theutterance \stepcounter{utterance}  
    & & & \multicolumn{4}{p{0.6\linewidth}}{
        \cellcolor[rgb]{0.9,0.9,0.9}{
            \makecell[{{p{\linewidth}}}]{
                \texttt{\tiny{[P1$\langle$GM]}}
                \texttt{Game Violation: ParseError: Response is not valid JSON.} \\
\texttt{Output ONLY:} \\
\texttt{One JSON with confidence\_score, current\_mode, input\_rationale, and hypothesis keys ...} \\
\texttt{OR} \\
\texttt{SOLVE: ```python ... ```} \\
            }
        }
    }
     \\ \\

    \theutterance \stepcounter{utterance}  
    & \multicolumn{4}{p{0.6\linewidth}}{
        \cellcolor[rgb]{0.561,0.737,0.561}{
            \makecell[{{p{\linewidth}}}]{
                \texttt{\tiny{[P1$\rangle$GM]}}
                \texttt{\{"current\_mode": "probing", "input\_rationale": "I have three examples: f(99, 1000) = 9, f({-}99, 0) = 9, and f(2, {-}2) = 4. The second argument seems irrelevant. Let's look at the first argument. 99 {-}> 9. {-}99 {-}> 9. 2 {-}> 4. This is strange. 2 squared is 4. 99? 9+9=18, 1+8=9. Maybe it's the digital root? Digital root of 99 is 9. Digital root of {-}99 is 9. Digital root of 2 is 2, not 4. So not digital root. Maybe it's the square of the first digit? For 99, first digit is 9, 9\^{}2=81, no. Maybe it's the last digit squared? 9\^{}2=81, no. Let's look at 2 {-}> 4. 2\^{}2=4. What about 99? If we take the last digit 9, 9\^{}2=81, no. Maybe it's the number of digits squared? 99 has 2 digits, 2\^{}2=4, no. Maybe it's the sum of digits squared? (9+9)\^{}2=} \\
            }
        }
    }
    & & \\ \\

    \theutterance \stepcounter{utterance}  
    & & & \multicolumn{4}{p{0.6\linewidth}}{
        \cellcolor[rgb]{0.9,0.9,0.9}{
            \makecell[{{p{\linewidth}}}]{
                \texttt{\tiny{[P1$\langle$GM]}}
                \texttt{Game Violation: ParseError: Response is not valid JSON.} \\
\texttt{Output ONLY:} \\
\texttt{One JSON with confidence\_score, current\_mode, input\_rationale, and hypothesis keys ...} \\
\texttt{OR} \\
\texttt{SOLVE: ```python ... ```} \\
            }
        }
    }
     \\ \\

    \theutterance \stepcounter{utterance}  
    & \multicolumn{4}{p{0.6\linewidth}}{
        \cellcolor[rgb]{0.561,0.737,0.561}{
            \makecell[{{p{\linewidth}}}]{
                \texttt{\tiny{[P1$\rangle$GM]}}
                \texttt{\{"current\_mode": "probing", "input\_rationale": "I have three examples: f(99, 1000) = 9, f({-}99, 0) = 9, and f(2, {-}2) = 4. The second argument seems irrelevant. Let's look at the first argument. 99 {-}> 9. {-}99 {-}> 9. 2 {-}> 4. This is strange. 2 squared is 4. 99? 9+9=18, 1+8=9. Maybe it's the digital root? Digital root of 99 is 9. Digital root of {-}99 is 9. Digital root of 2 is 2, not 4. So not digital root. Maybe it's the square of the first digit? For 99, first digit is 9, 9\^{}2=81, no. Maybe it's the last digit squared? 9\^{}2=81, no. Let's look at 2 {-}> 4. 2\^{}2=4. What about 99? If we take the last digit 9, 9\^{}2=81, no. Maybe it's the number of digits squared? 99 has 2 digits, 2\^{}2=4, no. Maybe it's the sum of digits squared? (9+9)\^{}2=} \\
            }
        }
    }
    & & \\ \\

    \theutterance \stepcounter{utterance}  
    & & & \multicolumn{2}{p{0.3\linewidth}}{
        \cellcolor[rgb]{0.9,0.9,0.9}{
            \makecell[{{p{\linewidth}}}]{
                \texttt{\tiny{[GM$|$GM]}}
                \texttt{Game Verdict: LOSS (turn limit reached)} \\
            }
        }
    }
    & & \\ \\

\end{supertabular}
}

\end{document}
